\section{零中心ファイバーのCayley三次曲面と楕円境界}
\label{sec:quotient-cayley}
\begin{proposition}[選ばれた反標準因子を持つCayley表示]
$C_0=Z(D_0)$の射影閉包$\overline S\subset\mathbb P^3_{U,V,W,T}$は
\[
 P=49U^3-14U^2T+UT^2+(11U-T)VW+V^3-W^3=0
\]
で定まり、Cayley三次曲面である。
$\omega^2+\omega+1=0$として
\[
 \begin{pmatrix}U\\V\\W\\T\end{pmatrix}
 =\begin{pmatrix}
 -9&-1&-1&-1\\
 0&-4&-4\omega&-4\omega^2\\
 0&4&4\omega^2&4\omega\\
 -63&1&1&1
 \end{pmatrix}
 \begin{pmatrix}a_0\\a_1\\a_2\\a_3\end{pmatrix}
\]
と置くと、変換は可逆で
\[
 P=-1728\sum_{i<j<k}a_i a_j a_k,\qquad
 T=-63a_0+a_1+a_2+a_3.
\]
$\operatorname{Spec}C_0=\overline S\setminus E$で、
$E=\overline S\cap\{T=0\}$は滑らかな反標準楕円曲線である。
\end{proposition}
\begin{proof}
変換行列の行列式は$-3456(2\omega+1)\ne0$。
$\omega^2+\omega+1$で剰余を取る直接代入で二式を得る。
無限遠の曲線は
\[
 E:\quad49U^3+11UVW+V^3-W^3=0.
\]
その三偏微分が零なら、いずれかの座標が零のとき全座標が零となる。
全座標が非零なら$147U^2=-11VW$、$3V^2=-11UW$、$-3W^2=-11UV$を掛け、
$1323=1331$という矛盾を得る。
従って$E$は滑らかで、無限遠に曲面の特異点もない。
三次超曲面の随伴公式$\omega_{\overline S}=\mathcal O_{\overline S}(-1)$から反標準性が従う。
\end{proof}

\begin{corollary}[Brauer指数6のtwisted NCCR]
$D_0$は正規Gorenstein曲面$C_0$のtwisted noncommutative crepant resolutionである。
そのBrauer指数は6であり、通常の$\operatorname{End}_{C_0}(N)$型のNCCRではない。
\end{corollary}
\begin{proof}
$C_0$は四つの$A_1$点以外で滑らかで、超曲面なので正規Gorenstein。
$R_1$の中心Macaulay性を中心正則元$g$で割ると$D_0$は$C_0$上MCMである。
従って反射的で、高さ1ではAzumaya、かつ大域次元2である。
\cite[Definitions~2.1, 2.2]{VdBOverview}のtwisted NCCRの条件を満たす。
一般中心分数体上は次数6の斜体なので、通常の自己準同型環なら行列環となって矛盾する。
局所完備化での非捩れ模型と、この大域Brauer障害を区別する。
\end{proof}
\begin{remark}[非可換affine del Pezzoとの比較]
この中心曲面は、選んだ滑らかな反標準楕円曲線を除いたCayley三次曲面であり、
\cite[\S1.1, Theorems~3.4.4, 3.4.5]{EGDelPezzo}の非可換affine del Pezzoの具体的な有限PI例と比較できる。
同文献のgenericな一生成元中心を、この特殊PI値の全中心へ適用してはいけない。
また定数項変形のPBW・CY性自体は\cite[Theorem~3.3.2]{EGDelPezzo}の既知の一般定理の特殊化であり、
本稿では具体的なsyzygyと自己双対解消も与えた。
Brauer類の具体的なsymbol表示、大域stack上の実現と大域対象の復元は未解決である。
\end{remark}

\section{古典的pencil、量子的判別式と非被約な点加群}
\label{sec:quotient-classical-comparison}
\begin{proposition}[古典的pencilと三様体の判別式]
\[
 F_\alpha=\frac{x^3+y^3+z^3}{3}+\alpha uvx+u^3+v^3
\]
とし、$v=au$の曲面切断を考える。
有限pencilパラメータでの特異条件と三様体の特異条件はそれぞれ
\[
 \Delta_\alpha(a)=4\alpha^3a^3+9(1+a^3)^2=0,\qquad
 F_\alpha\text{が特異}\ \Longleftrightarrow\ \alpha^3=-9.
\]
$p(x)=x^3/3+\alpha ax+1+a^3$に対し
$\operatorname{Res}_x(p,p')=\Delta_\alpha(a)/9$で、無限遠切断は滑らかである。
さらに
\[
 \operatorname{Disc}_a\Delta_\alpha
 =2^{12}3^{14}\alpha^9(\alpha^3+9)^3,\qquad
 \Delta_\alpha(\omega)=4(\alpha^3+9)\quad(\omega^3=1).
\]
\end{proposition}
\begin{proof}
$u=0$上のFermat平面三次曲線は滑らかであり、$u=1$の特異点では$y=z=0$。
従って曲面では$p,p'$の共通根を調べればよい。
三様体の勾配は
$y=z=0$、$x^2+\alpha uv=0$、$\alpha vx+3u^2=0$、$\alpha ux+3v^2=0$を与える。
射影点では$u,v\ne0$で、$u=1$の消去は
$3x-\alpha^2v^2=v^3-1=\alpha^3+9=0$となる。
従って三特異点は$[\alpha^2\omega^2/3:0:0:1:\omega]$である。
残りは一変数resultantとdiscriminantの厳密計算から従う。
\end{proof}
\begin{remark}[重複度を含むresultantの規格化]
$\operatorname{Res}(X_1^2,\ldots,X_n^2)=1$で規格化すると
\[
 \operatorname{Res}(\nabla F_\alpha)=3^8(\alpha^3+9)^{12},\qquad
 \operatorname{Res}(\nabla(F_\alpha|_{v=au}))=\Delta_\alpha(a)^4.
\]
最初は$(x,u,v)$の三変数resultant $9(\alpha^3+9)^3$の四乗である。
$x=1$における$\CC[u,v]/(3u^2+\alpha v,3v^2+\alpha u)$の基底$1,u,v,uv$を使うと、
$1+\alpha uv$の乗法行列式は$(\alpha^3+9)^3/729$となる。
最高次係数$3^8$とPoisson積公式\cite[equation~(2)]{DKSResultant}で三変数の式を得る。
切断の式も、残る二変数の勾配resultantと$y^2,z^2$の二つの平方から従う。
これらは古典的判別式であり、量子的なtails判別式ではない。
\end{remark}

\begin{proposition}[正則性を区別しない定常点の粗い局所データ]
$h\ne0$、$\alpha\ne0$、$\alpha^3\ne-9$とする。
$A_{\alpha,h}$の中心一次元が非零となるpoint moduleのスキームは、
六点を台とし、各局所環が$\CC[y,z]/(y^2,z^2)$である長さ24の有限スキームとなる。
\end{proposition}
\begin{proof}
可換試験環上で$u$が可逆なら、中心性の多重線型関係
$u_i a_{i+1}=a_i u_{i+1}$から$u_i=1$に規格化した点列は一定となる。
従ってこのchartではpoint moduleは可換化で表される。
可換化の関係は$y^2=z^2=0$、$x^2+\alpha v=0$である。
$v=-x^2/\alpha$を消去すると中心三次式は
\[
 1-\frac{2x^3}{3}-\frac{x^6}{\alpha^3}-\frac{2h^2}{3}xyz=0.
\]
$y=z=0$での$x^3$の方程式は
$3q^2+2\alpha^3q-3\alpha^3=0$で、判別式は$4\alpha^3(\alpha^3+9)$。
二根は相異なり非零なので六つの$x$を与える。
$x$方向の微分は各点で単元であり、形式的陰関数定理で$x$を$yz$の式として消せる。
従って各局所環は表示した長さ4の環である。
$u=0,v\ne0$の可換化では$x^2=y^2=z^2=0$かつ$v^3=0$となり矛盾するので、
中心座標が非零の点は全てこのchartに含まれる。
\end{proof}
\begin{remark}
$h=1$で$\alpha^3=-243/28$なら、古典三様体と楕円境界は滑らかでもtailsは非正則である。
$\alpha=1$の正則例とこの三つの非正則例は、上の定常点の台の個数、
局所環型と全長が一致する。埋込みを含む全point schemeの同型を主張してはいない。
$\alpha\ne0,\alpha^3\ne-9$の古典pencilには六つの特異曲面があり、
$\alpha=0$では判別式は$9(a^3+1)^2$、三切断は滑らかなFermat曲線上の錐である。
前者の各特異曲面は、$x$方向の非退化二次項を消去すると
局所Jacobian環$\CC[y,z]/(y^2,z^2)$（長さ4）を持つ。
$\alpha=0$の三つの錐の頂点では$\CC[x,y,z]/(x^2,y^2,z^2)$となり長さ8である。
全切断の滑らかさやpencil判別式の重根の有無を、三様体自体の滑らかさと混同しない。
\end{remark}
